\chapter{Fever in Children}

\section*{Definition}
Fever is an elevated body temperature caused most often by infection. In children, age, appearance, breathing, hydration, and behavior are more important than the number alone.

\section*{When to See a Doctor}
Emergency assessment is needed for a baby younger than three months with a temperature of 38°C or higher, breathing difficulty, blue color, seizure, stiff neck, non-blanching rash, severe dehydration, unusual sleepiness, or a child who looks seriously ill. Arrange prompt evaluation for persistent, recurrent, unexplained, or function-limiting symptoms.

\section*{How It Develops}
Fever is a regulated rise in body temperature caused by immune signals, most often during infection. The child's appearance, breathing, circulation, hydration, and age are more useful than temperature alone in judging severity.

\section*{Causes}
\begin{itemize}
  \item Viral or bacterial infection, immunization response, inflammatory disease, and less commonly serious bloodstream, lung, urinary, brain, or bone infection.
  \item Medication effects, environmental exposures, and underlying chronic disease may also contribute.
  \item Serious causes are less common but must be considered when symptoms are sudden, progressive, severe, or associated with systemic illness.
\end{itemize}

\section*{Related Symptoms}
\begin{itemize}
  \item Pain, fever, fatigue, nausea, weakness, or changes in normal function
  \item Bleeding, swelling, discharge, breathing difficulty, or neurologic symptoms when a serious cause is present
\end{itemize}

\section*{Medical Evaluation}
\begin{itemize}
  \item History addresses onset, duration, severity, triggers, injuries, exposures, medications, medical conditions, age, and pregnancy status when relevant.
  \item Examination includes vital signs and focused assessment of the relevant organ system, neurologic status, hydration, and function.
  \item Testing may include blood or urine studies, cultures, imaging, physiologic testing, fetal or pediatric assessment, or specialist examination.
\end{itemize}

\section*{Treatment Options}
Offer fluids and comfort measures. Paracetamol or ibuprofen may be used at age-appropriate doses when the child is distressed, but antipyretics do not prevent febrile seizures; antibiotics are used only for selected bacterial infections.

\section*{Self-Care and Home Management}
Avoid known triggers, protect the affected area, maintain appropriate hydration, and record timing and associated symptoms. Use only age- or pregnancy-appropriate medicines recommended by a clinician. Do not delay emergency care while trying home treatment.

\section*{References}
\begin{thebibliography}{99}
\bibitem{fever_in_children:nice} National Institute for Health and Care Excellence. Fever in under 5s: assessment and initial management. NICE guideline NG143; 2019 (updated 2021).
\bibitem{fever_in_children:aap} Pantell RH, Roberts KB, Adams WG, et al. Evaluation and management of well-appearing febrile infants 8 to 60 days old. \textit{Pediatrics}. 2021;148:e2021052228.
\bibitem{fever_in_children:antipyretics} Sullivan JE, Farrar HC. Fever and antipyretic use in children. \textit{Pediatrics}. 2011;127:580--587.
\end{thebibliography}
